\documentclass[12pt,english,a4paper]{article}

\usepackage[a4paper, total={6in, 9in}]{geometry}
\usepackage[english]{babel}
\usepackage{graphicx} % Required for inserting images
\usepackage{url}
\usepackage{hyperref}
\usepackage{csquotes}
%\usepackage{cite}
%\usepackage{times}
\usepackage[backend=biber, sorting=none, maxnames=10,minnames=10]{biblatex} %bibtex or biber
\addbibresource{literature.bib}

\usepackage{fancyhdr}
\fancypagestyle{plain}{
    \fancyhf{}
    \fancyfoot[C]{\thepage}
    \renewcommand{\thepage}{\ifnum\value{page}<10 0\fi\arabic{page}}
    \renewcommand{\headrulewidth}{0pt} % Remove header line
}

\pagestyle{plain}

\title{Deep Learning in Video Streaming Services\thanks{Semester project in the subject Engineering Methods, academic year 2024/2025, supervisor: Pavol Baťalík}}

\author{Arseniy Malyuk\\[2pt]
	{\small Slovak University of Technology in Bratislava}\\
	{\small Faculty of Informatics and Information Technologies}\\
	{\small \texttt{xmalyuk@stuba.sk}}
	}
\date{\small November 2024}

\begin{document}

\maketitle

\begin{abstract}

    The purpose of this article is to explore how video streaming services such as Netflix, Amazon Prime Video, Disney+, etc. implement deep learning for personalized recommendations. The main area of focus here will be evaluation of deep learning approaches for analyzing large quantities of user data such as activity logs and preferences. We will explore how such algorithms are used for generation of accurate recommendations, increasing user satisfaction or even improving viewing rates. We will also consider using NLP and video content analysis methods for recommendation improvement. In particular, we would like to know how various streaming services operate and implement deep learning for assessing users’ interaction activities with the content such as views, ratings, comments, etc. This will allow us to see the personalized recommendations that create based on users’ interests and the information they already have. Towards this, we will look at some practical aspects of the deep learning application in recommender systems and evaluate the performance. This will guide us on how best such solutions can be put to use in practice and the possible limitations that might be encountered. At the end, we will discuss how improvement in these methods affects the quality and effectiveness of the streaming services. The main source of inspiration for this topic was the article "An overview of video recommender systems: state-of-the-art and research issues" \cite{Lubos2023} by Sebastian Lubos, Alexander Felfernig, and Markus Tautschnig. This article provides a comprehensive overview of current algorithms and methods in the field of video recommendation.
    
\end{abstract}

\newpage

\section{Introduction}

\hspace*{2em}Each year, television is becoming less and less popular and the initiative of content consumption is shifting to video streaming services such as Netflix, Amazon Prime Video, Disney+, etc., where you do not need to switch between channels to find what you like, but simply go to the "Recommended" section. But how exactly do recommendations work and how can they "predict" user preferences? This seems to be a magical ability, but it is actually the result of sophisticated algorithms and deep learning models that are able to analyse and learn from massive amounts of data. It is thanks to deep learning that these systems are able to accurately predict what users would like and provide unique content for each of them.

Deep learning opens up opportunities to categorise users based on their previous activity (likes, comments, shares, etc.). By detecting certain patterns, machine learning models can predict the type of content that a user is likely to view in the future.

However, the history of user activity is not the only thing that forms the basis for final recommendations. Natural language processing (NLP) and video analysis are other important tools for creating personalised recommendations. NLP is used specifically to analyse the text itself, which can be found in comments, reviews, search terms, etc. This is done to help the system better understand emotions and context. On the other hand, video analysis includes the study of genres, themes, and even visual elements, which helps to better understand user preferences.

Therefore, recommendation systems based on deep learning are an irreplaceable tool when it comes to analysing and predicting user preferences. However, the principle of deep learning in video streaming services is too extensive to be described in just a few sentences. For this reason, in next sections we will take a closer look at how deep learning is used in video streaming services.

\section{Types of Recommendation Systems}

\hspace*{2em}Streaming platforms use various types of recommendation systems, each designed to capture different aspects of user preference. The major categories include content-based systems, collaborative filtering, hybrid approaches and group recommendations.

\subsection{Content-Based Recommendation}

\hspace*{2em}In content-based systems, recommendations are generated by analyzing features associated with the content itself, such as genre, cast, theme, etc. For example, if a viewer often watches historical dramas, the system will recommend similar content. These systems create user profiles based on the metadata of previously watched shows or movies, so suggestions are based on specific content attributes rather than other users’ choices \cite{Lubos2023,Peng2024}.

\subsection{Collaborative Filtering}

\hspace*{2em}Collaborative filtering suggests content, based on what similar users have watched. There are two key approaches here:

\begin{itemize}
    \item \textbf{User-Based Collaborative Filtering:} This method recommends content to a user, based on the preferences of others with similar viewing patterns. For instance, if User A and User B share a taste in thrillers, User A’s preferred shows might be suggested to User B \cite{Lubos2023,Dang2021}.
    
    \item \textbf{Item-Based Collaborative Filtering:} Instead of focusing on user similarity, item-based filtering finds similarities among shows or movies themselves. If a particular group of users has engaged with both "Game of Thrones" and "The Witcher," item-based filtering might recommend one to someone who watched the other \cite{Lubos2023,Dang2021}.
\end{itemize}

Collaborative filtering, which leverages the preferences of a large user base, is highly effective in personalizing recommendations. Netflix, for instance, has successfully implemented this method to create precise recommendations for its subscribers \cite{Steck20217,Mu2023}.

\subsection{Hybrid Recommendation Systems}

\hspace*{2em}Hybrid recommendation systems combine aspects of content-based and collaborative filtering to increase accuracy and relevance. By drawing on both individual and broader user base behaviors, hybrid systems address potential shortcomings of each method. Netflix’s model, which uses both approaches, is a strong example of how hybrid systems can effectively balance variety and specificity in recommendations\\\cite{Lubos2023,Dang2021}.

\subsection{Group Recommendation Systems}

\hspace*{2em}Group recommendations are tailored for users who watch content together. These systems aggregate individual preferences to recommend content that appeals to multiple people in a shared viewing session. For example, family or friend groups can receive suggestions based on everyone’s preferences, creating a viewing experience that satisfies the group as a whole \cite{Lubos2023}.

\section{Implementation of Deep Learning}

\subsection{Neural Networks and Algorithms}

\hspace*{2em}Deep learning models, particularly Convolutional Neural Networks (CNNs) and Recurrent Neural Networks (RNNs) are instrumental in powering recommendation systems used by streaming services.

\begin{itemize}
    \item \textbf{CNNs} are especially effective at processing visual and video data, enabling the extraction of relevant features from images and scenes. This capability is vital for understanding the visual elements of content, which can inform more accurate recommendations \cite{Hande2023,Alzubaidi2021}.

    \item \textbf{RNNs} are adept at handling sequential data, such as the patterns in user viewing habits over time. By analyzing how a viewer's preferences evolve, RNNs contribute to the creation of dynamic user profiles that reflect shifting tastes \cite{Alzubaidi2021}.

\end{itemize}
Netflix, for example, uses variational autoencoders and multi-head attention mechanisms to dynamically adapt recommendations based on context, such as the user’s device or time of day \cite{Gomez-Uribe2015,Jayalakshmi2022}.

\subsection{Netflix’s Recommendation System}

\hspace*{2em}Netflix’s recommendation system is one of the most advanced in the industry, blending various deep learning techniques to deliver highly personalized content. By analyzing user data in real time, Netflix’s algorithms adjust recommendations based on recent behavior, providing contextually relevant content that changes with individual viewing habits. This adaptability is a core strength of deep learning in the streaming industry, allowing for real-time, personalized engagement \cite{Steck20217,Gomez-Uribe2015}.

\section{Evaluation of Deep Learning Approaches}

\hspace*{2em}Evaluating deep learning models is essential to ensure they are meeting goals for relevance, engagement and efficiency. Key performance indicators include accuracy and user satisfaction.

\subsection{Accuracy of Recommendations}

\hspace*{2em}To assess the accuracy of recommendations, streaming platforms typically rely on precision, recall, and F1 scores:

\begin{itemize}
    \item \textbf{Precision} measures the proportion of recommended content that users engage with, providing insight into the relevance of the suggestions \cite{Fayyaz20201}.
    
    \item \textbf{Recall} assesses the effectiveness of the system in delivering all relevant content to users, evaluating how well the algorithm captures user interests \cite{Fayyaz20201}.
    
    \item The \textbf{F1 score} offers a comprehensive assessment by combining precision and recall into a single metric, allowing for a balanced evaluation of recommendation quality \cite{Fayyaz20201}.
\end{itemize}

\subsection{User Satisfaction}

\hspace*{2em}To evaluate user satisfaction, platforms track engagement metrics such as retention rates and content engagement. These metrics help streaming platforms fine-tune their models, ensuring that recommendations remain relevant and meaningful over time. User satisfaction data is invaluable for updating algorithms to reflect changing viewer preferences \cite{Souane2023}.

\section{NLP and Video Content Analysis}

\subsection{NLP Techniques}

\hspace*{2em}Natural Language Processing (NLP) interprets user-generated text, such as reviews, comments, search terms etc. Using sentiment analysis and topic modeling, NLP helps systems understand not only what users are watching, but how they feel about it. Sentiment analysis classifies text as positive, neutral or negative, while topic modeling categorizes comments by theme, allowing platforms to refine recommendations, based on emotional tone and specific interests \cite{Hande2023,Saini202312347}.

\subsection{Video Content Analysis}

\hspace*{2em}Video content analysis examines visual and thematic elements in video content, from genre to specific visual features. By analyzing patterns within video data, platforms can enhance recommendations with detailed content information, creating a refined viewing experience that aligns more closely with user preferences. For example, a viewer interested in "space" themes might be recommended other visually similar content \cite{Souane2023,Saini202312347}.

\section{Practical Applications and Performance}

\section{Conclusion}

\newpage

\printbibliography

\end{document}